% =====================================================================
% OBJETIVOS DE DESARROLLO SOSTENIBLE
% =====================================================================
% BORRADOR. La normativa pide al menos 200 palabras; este texto tiene unas 460.
% Revisa que te representa, y si tu tutor sugiere otro ODS, cambia el enfoque.
%
% Los subtitulos van con \noindent\textbf{} y NO con \paragraph{}, porque memoir
% mete los \paragraph en el indice general y este es un apartado de preliminares
% que no debe aportar entradas.
%
% ESTILO: sin negritas en el texto corrido y sin incisos entre rayas.
% =====================================================================

Este trabajo contribuye principalmente al ODS 9, Industria, innovación e
infraestructura, y de forma secundaria al ODS 16, Paz, justicia e instituciones
sólidas.

\medskip

\noindent\textbf{ODS 9. Infraestructura resiliente.} La meta 9.1 pide
desarrollar infraestructuras fiables y resilientes, y la 9.c, aumentar el acceso
a las tecnologías de la información. Las redes de comunicaciones son hoy la
infraestructura sobre la que se apoyan la sanidad, la energía, el transporte y
la administración pública, y su fiabilidad depende de que los ataques se
detecten a tiempo. Un sistema de detección de intrusiones es, en ese sentido,
infraestructura que protege infraestructura.

La contribución concreta de esta memoria no es un detector mejor, sino algo que
condiciona a todos: mostrar que la forma habitual de evaluarlos no mide lo que
se cree. Se demuestra que la evaluación dentro de un mismo día está saturada y
no distingue entre representaciones, y que buena parte del éxito aparente del
análisis de contenido es reconocimiento de la herramienta con que se generó el
ataque, no del ataque en sí. Desplegar un sistema validado con ese protocolo
produce una falsa sensación de seguridad, que es peor que no tener sistema,
porque se confía en una defensa que no defiende. Aportar métodos de
verificación, como una \emph{ground-truth} comprobada sobre los datos, pruebas
de evasión y medición de la transferencia entre dominios, es contribuir a que la
infraestructura digital sea resiliente de verdad y no solo sobre el papel.

Hay además una contribución en eficiencia que la meta 9.4 recoge al pedir
industrias sostenibles con uso eficiente de los recursos. El trabajo mide que
las arquitecturas más costosas contrastadas no mejoran los resultados y llegan a
consumir entre 2 y 49 veces más cómputo. Documentar cuándo un modelo grande no
aporta nada evita gasto energético inútil, en un campo donde la tendencia por
defecto es escalar el modelo.

\medskip

\noindent\textbf{ODS 16. Instituciones sólidas.} La meta 16.4 aborda la
reducción de las corrientes financieras y de armas ilícitas, y la 16.6, el
desarrollo de instituciones eficaces y transparentes. El cibercrimen es hoy uno
de los vectores principales de actividad ilícita transfronteriza, y su coste se
estima en billones de euros anuales. Mejorar la capacidad de detección de
intrusiones, y sobre todo la capacidad de saber cuándo esa detección es real y
cuándo es un artefacto de la medición, apoya directamente la protección de las
instituciones públicas y privadas frente a esa actividad.

La dimensión de transparencia también aplica. Este trabajo publica su
metodología, sus datos derivados y sus errores, incluidos los que tuvo que
corregir sobre la marcha. La investigación en seguridad que solo reporta sus
aciertos contribuye a un cuerpo de conocimiento poco fiable; documentar los
resultados negativos y las hipótesis refutadas es una forma modesta pero
concreta de sostener la meta 16.6.
